% =============================================================
% preamble_manual.tex — оформление учебно-методического пособия
% Основа — latex/preamble.tex (СТП БГУИР 01-2024): A4, поля 30/15/20/20 мм,
% Times 14 pt, одинарный интервал, абзацный отступ 1,25 см,
% номер страницы внизу справа, заголовки разделов с новой страницы,
% подписи «Рисунок 1.1 – ...» под рисунком и «Таблица 1.1 – ...» над таблицей.
% =============================================================

\usepackage[utf8]{inputenc}
\usepackage[T2A]{fontenc}
\usepackage[russian]{babel}

% Шрифт Times (текст и формулы)
\usepackage{amsmath}
\usepackage{tempora}
\usepackage[varg]{newtxmath}
\usepackage{icomma}          % десятичная запятая в формулах: $0,5$
\usepackage{mathtools}

% Поля и интервалы
\usepackage{geometry}
\geometry{left=30mm, right=15mm, top=20mm, bottom=20mm, footskip=10mm}
\usepackage{setspace}
\singlespacing
\usepackage{indentfirst}
\setlength{\parindent}{1.25cm}
\setlength{\parskip}{0pt}
\pretolerance=1000
\tolerance=2000
\emergencystretch=3em
\clubpenalty=10000
\widowpenalty=10000

% Нумерация страниц: внизу справа
\usepackage{fancyhdr}
\pagestyle{fancy}
\fancyhf{}
\fancyfoot[R]{\thepage}
\renewcommand{\headrulewidth}{0pt}
\fancypagestyle{plain}{\fancyhf{}\fancyfoot[R]{\thepage}\renewcommand{\headrulewidth}{0pt}}

% Графика, таблицы
\usepackage{graphicx}
\usepackage{float}
\usepackage[table]{xcolor}
\usepackage{booktabs}
\usepackage{tabularx}
\usepackage{multirow}
\usepackage{longtable}
\usepackage{array}
\usepackage{makecell}
\usepackage{needspace}
\usepackage{tikz}
\usetikzlibrary{arrows.meta, positioning, calc, shapes.geometric, shapes.misc, fit, backgrounds, matrix}

\definecolor{cblue}{HTML}{2A78D6}
\definecolor{corange}{HTML}{EB6834}
\definecolor{cgreen}{HTML}{1BAF7A}
\definecolor{cyellow}{HTML}{EDA100}
\definecolor{cgray}{HTML}{898781}
\definecolor{clight}{HTML}{E1E0D9}

% Подписи
\usepackage{caption}
\captionsetup[figure]{labelsep=endash, justification=centering, font=normalsize, name=Рисунок, skip=4pt}
\captionsetup[table]{labelsep=endash, justification=raggedright, singlelinecheck=false,
  font=normalsize, name=Таблица, skip=3pt}

% Плавающие объекты
\renewcommand{\topfraction}{0.9}
\renewcommand{\bottomfraction}{0.8}
\renewcommand{\textfraction}{0.07}
\renewcommand{\floatpagefraction}{0.85}
\setcounter{topnumber}{3}
\setcounter{bottomnumber}{2}
\setcounter{totalnumber}{4}
\setlength{\textfloatsep}{10pt plus 3pt minus 3pt}
\setlength{\floatsep}{8pt plus 2pt minus 2pt}
\setlength{\intextsep}{8pt plus 2pt minus 2pt}
\setlength{\abovedisplayskip}{6pt plus 2pt minus 2pt}
\setlength{\belowdisplayskip}{6pt plus 2pt minus 2pt}
\setlength{\abovedisplayshortskip}{4pt}
\setlength{\belowdisplayshortskip}{4pt}

% Заголовки: раздел с новой страницы, полужирный, прописные, с абзацного отступа
\usepackage{titlesec}
\newif\ifnosectionbreak
\titleformat{\section}
  {\ifnosectionbreak\global\nosectionbreakfalse\else\clearpage\fi\normalfont\fontsize{14}{17}\selectfont\bfseries\raggedright}
  {\thesection}{0.5em}{\MakeUppercase}
\titlespacing*{\section}{1.25cm}{0pt}{12pt}
\titleformat{\subsection}
  {\normalfont\fontsize{14}{17}\selectfont\bfseries\raggedright}{\thesubsection}{0.5em}{}
\titlespacing*{\subsection}{1.25cm}{12pt}{8pt}
\titleformat{\subsubsection}
  {\normalfont\fontsize{14}{17}\selectfont\bfseries\raggedright}{\thesubsubsection}{0.5em}{}
\titlespacing*{\subsubsection}{1.25cm}{10pt}{6pt}

% Заголовок без номера по центру (введение, список источников)
\newcommand{\sectioncenter}[1]{%
  \clearpage\phantomsection\addcontentsline{toc}{section}{#1}%
  \begin{center}\textbf{\MakeUppercase{#1}}\end{center}\vspace{6pt}}

% Название части (теоретическая часть, практикум): следующий раздел идёт на той же странице
\newcommand{\parttitle}[1]{%
  \clearpage\phantomsection\addcontentsline{toc}{section}{\MakeUppercase{#1}}%
  \begin{center}\fontsize{16}{20}\selectfont\bfseries\MakeUppercase{#1}\end{center}\vspace{8pt}%
  \global\nosectionbreaktrue}

% Лабораторная работа: номер раздела «Л1», рисунки «Л1.1»
\newcommand{\labsection}[2]{%
  \clearpage\refstepcounter{section}\phantomsection
  \addcontentsline{toc}{section}{Лабораторная работа №\arabic{section}. #2}%
  \begin{center}\bfseries ЛАБОРАТОРНАЯ РАБОТА №\arabic{section}\\[2pt]\MakeUppercase{#2}\end{center}\vspace{4pt}}
% Рубрика внутри лабораторной работы
\newcommand{\labpart}[1]{\par\addvspace{6pt}\noindent\hspace{1.25cm}\textbf{#1}\par\nobreak\vspace{3pt}\nobreak}

% Приложение
\newcommand{\appendixsection}[3]{%
  \clearpage\refstepcounter{section}\phantomsection
  \addcontentsline{toc}{section}{Приложение \thesection. #3}%
  \begin{center}\textbf{ПРИЛОЖЕНИЕ \thesection}\\(#2)\\[2pt]\textbf{#3}\end{center}\vspace{6pt}}

% Содержание
\usepackage{tocloft}
\renewcommand{\cftsecfont}{\normalfont}
\renewcommand{\cftsecpagefont}{\normalfont}
\renewcommand{\cftsubsecfont}{\normalfont}
\renewcommand{\cftsubsecpagefont}{\normalfont}
\renewcommand{\cftsecleader}{\cftdotfill{\cftdotsep}}
\renewcommand{\cftdotsep}{1}
\setlength{\cftbeforesecskip}{2pt}
\setlength{\cftsecnumwidth}{1.6em}
\setlength{\cftsubsecindent}{1.6em}
\setlength{\cftsubsecnumwidth}{2.3em}
\addto\captionsrussian{\renewcommand{\contentsname}{\hfill\textbf{СОДЕРЖАНИЕ}\hfill}}
\renewcommand{\cftaftertoctitle}{\hfill}
\setcounter{tocdepth}{2}

% Списки: тире, без лишних интервалов
\usepackage{enumitem}
\setlist{beginpenalty=10000}   % не отрывать список от заголовка рубрики или вводной строки
\setlist[itemize]{noitemsep, topsep=2pt, leftmargin=*, labelindent=1.25cm, label=--}
\setlist[enumerate]{noitemsep, topsep=2pt, leftmargin=*, labelindent=1.25cm, label=\arabic*)}
\newlist{steps}{enumerate}{1}
\setlist[steps]{noitemsep, topsep=2pt, leftmargin=0pt, itemindent=1.25cm+1.5em, labelwidth=1.3em,
  labelsep=0.4em, label=\arabic*.}
\newlist{questions}{enumerate}{1}
\setlist[questions]{noitemsep, topsep=2pt, leftmargin=0pt, itemindent=1.25cm+1.5em, labelwidth=1.3em,
  labelsep=0.4em, label=\arabic*.}

% Листинги
\usepackage{listings}
\lstset{
  language=Python,
  basicstyle=\ttfamily\fontsize{9}{10.8}\selectfont,
  keywordstyle=\bfseries,
  commentstyle=\itshape\color{cgray},
  frame=single, rulecolor=\color{clight},
  breaklines=true, showstringspaces=false,
  columns=fullflexible, keepspaces=true,
  xleftmargin=2pt, xrightmargin=2pt,
  inputencoding=utf8, extendedchars=true,
  literate={а}{{\cyra}}1 {б}{{\cyrb}}1 {в}{{\cyrv}}1 {г}{{\cyrg}}1 {д}{{\cyrd}}1 {е}{{\cyre}}1
           {ё}{{\cyryo}}1 {ж}{{\cyrzh}}1 {з}{{\cyrz}}1 {и}{{\cyri}}1 {й}{{\cyrishrt}}1 {к}{{\cyrk}}1
           {л}{{\cyrl}}1 {м}{{\cyrm}}1 {н}{{\cyrn}}1 {о}{{\cyro}}1 {п}{{\cyrp}}1 {р}{{\cyrr}}1
           {с}{{\cyrs}}1 {т}{{\cyrt}}1 {у}{{\cyru}}1 {ф}{{\cyrf}}1 {х}{{\cyrh}}1 {ц}{{\cyrc}}1
           {ч}{{\cyrch}}1 {ш}{{\cyrsh}}1 {щ}{{\cyrshch}}1 {ъ}{{\cyrhrdsn}}1 {ы}{{\cyrery}}1
           {ь}{{\cyrsftsn}}1 {э}{{\cyrerev}}1 {ю}{{\cyryu}}1 {я}{{\cyrya}}1
           {А}{{\CYRA}}1 {Б}{{\CYRB}}1 {В}{{\CYRV}}1 {Г}{{\CYRG}}1 {Д}{{\CYRD}}1 {Е}{{\CYRE}}1
           {Ж}{{\CYRZH}}1 {З}{{\CYRZ}}1 {И}{{\CYRI}}1 {К}{{\CYRK}}1 {Л}{{\CYRL}}1 {М}{{\CYRM}}1
           {Н}{{\CYRN}}1 {О}{{\CYRO}}1 {П}{{\CYRP}}1 {Р}{{\CYRR}}1 {С}{{\CYRS}}1 {Т}{{\CYRT}}1
           {У}{{\CYRU}}1 {Ф}{{\CYRF}}1 {Х}{{\CYRH}}1 {Ц}{{\CYRC}}1 {Ч}{{\CYRCH}}1 {Ш}{{\CYRSH}}1
           {Э}{{\CYREREV}}1 {Ю}{{\CYRYU}}1 {Я}{{\CYRYA}}1 {—}{{\textemdash}}1 {→}{{$\to$}}1
}

% Ссылки
\usepackage{hyperref}
\hypersetup{colorlinks=true, linkcolor=black, urlcolor=black, citecolor=black,
  pdftitle={Генетические и эволюционные вычисления: нейронные сети и генетические алгоритмы},
  pdfauthor={А. В. Лебедевич, С. Н. Нестеренков}}
\usepackage{lastpage}

% Нумерация рисунков, таблиц и формул в пределах раздела
\numberwithin{figure}{section}
\numberwithin{table}{section}
\numberwithin{equation}{section}

% Список источников: «1 Автор...» с абзацного отступа (ГОСТ 7.1)
\makeatletter
\renewcommand{\@biblabel}[1]{#1}
\renewenvironment{thebibliography}[1]
  {\sectioncenter{Список использованных источников}%
   \list{\@biblabel{\@arabic\c@enumiv}}%
     {\settowidth\labelwidth{\@biblabel{#1}}%
      \setlength{\leftmargin}{0pt}\setlength{\itemindent}{1.25cm}%
      \addtolength{\itemindent}{\labelwidth}\addtolength{\itemindent}{\labelsep}%
      \setlength{\itemsep}{1pt}\setlength{\parsep}{0pt}%
      \usecounter{enumiv}\let\p@enumiv\@empty\renewcommand\theenumiv{\@arabic\c@enumiv}}%
   \sloppy}
  {\endlist}
\makeatother

% Удобные обозначения
\newcommand{\figdir}{figures}
\newcommand{\labfig}[2]{../lab_0#1/report/figures/#2}
\newcommand{\code}[1]{\texttt{#1}}
\newcommand{\bits}[1]{\texttt{#1}}
\DeclareMathOperator{\dec}{dec}
\DeclareMathOperator*{\argmin}{arg\,min}
\DeclareMathOperator*{\argmax}{arg\,max}

% Строка битов/генов в рамках (TikZ): \bitrow[x0]{y}{подпись}{0/white,1/cblue!20,...}
\newlength{\bitw}\setlength{\bitw}{4.4mm}
\newcommand{\bitrow}[4][0]{%
  \node[anchor=east, font=\small] at (#1, #2) {#3};
  \foreach \b/\c [count=\i] in {#4}{%
    \node[draw, minimum size=\bitw, inner sep=0pt, fill=\c, font=\ttfamily\small]
      at ($(#1, #2) + (\i*\bitw - 0.3\bitw, 0)$) {\b};}}
% Вертикальная линия разреза после гена k: \cutline[x0]{k}{y_верх}{y_низ}
\newcommand{\cutline}[4][0]{%
  \draw[corange, very thick, densely dashed] ($(#1, #3) + (#2\bitw + 0.2\bitw, 0.35)$) --
    ($(#1, #4) + (#2\bitw + 0.2\bitw, -0.35)$);}
